\begingroup
\emergencystretch=3em
\section{Every local heat coefficient and its exact residue receiver}
\label{sec:all-order-local-residue}

This section retains the original heat equation, spatial coordinate,
arithmetic coordinate, and full analytic unit. The primary heat source
is Brad Rodgers and Terence Tao,
\href{https://arxiv.org/abs/1801.05914v5}{\emph{The de Bruijn--Newman
constant is non-negative}}, original author TeX
\texttt{fmp-template.tex}, lines 78--101, equations \texttt{hoz},
\texttt{sas}, \texttt{phidef}, and \texttt{htdef}. Those author-source
lines were read directly; the full paper is not claimed reread here.
The preceding marked residue maps are
\href{https://github.com/KokunoYumeto/zeta-function-research-reader/blob/87e386271c488ce5914b9926c4d23d0857e0c367/workbenches/splitzero-tandem/continuations/20260923-boundary-residue-arithmetic/BOUNDARY_RESIDUE_ARITHMETIC_RECEIVERS.tex}
{RJ3--9 and RJ18--20 in the complete boundary/residue proof}.
The earlier fixed-test calculation is
\href{https://github.com/KokunoYumeto/zeta-function-research-reader/blob/87e386271c488ce5914b9926c4d23d0857e0c367/workbenches/splitzero-tandem/continuations/20260923-boundary-residue-arithmetic/RETAINED_FIXED_ARITHMETIC_HEAT_FLAGS.tex}
{HR3--6 in the retained fixed-arithmetic proof}.
Both complete sources accompany the calculation. Every map and coefficient
needed for the extension below is proved here.

\paragraph{LT1. The actual local family and the finite contour.}
Fix an actual real zero $x_0$ of order $m\ge1$ of the original
$H_0$, and write $x=Z-x_0$ without changing its scale. More generally
the proof applies to any jointly holomorphic solution of
$\partial_tH=-\partial_x^2H$ with
\[
 H_0(x_0+x)=x^mu(x),\qquad
 u(x)=\sum_{r\ge0}u_rx^r,\quad u_0\ne0,\qquad
 s=\rho+ix/2,\quad \rho=\tfrac12+ix_0/2.
 \tag{LT1}
\]
For the original family $H_0(Z)=\xi_{\rm R}(1/2+iZ/2)/8$.
Its joint holomorphy follows directly from the retained theta
integral: on $u\ge0$, its defining series has the bound
$|\Phi(u)|\le C e^{9u}\exp(-(\pi/2)e^{4u})$.
Indeed one factors $\exp(-(\pi/2)e^{4u})$ from each summand
and bounds the remaining convergent sum by a constant times
$\sum_{l\ge1}l^4\exp(-\pi(l^2-1/2))$.
On $|t|\le T$, $|Z|\le R$, the integrand and every fixed
parameter derivative are dominated by a polynomial in $u$
times $\exp(Tu^2+Ru+9u-(\pi/2)e^{4u})$, which is integrable.
Uniform convergence on compact parameter sets proves the
joint holomorphy used here, with the original kernel intact.
The test $\phi(x)$ is holomorphic on a fixed disk about $0$.
Choose a smaller disk on whose closure $u$ does not vanish and
whose boundary $\Gamma$ lies in the common holomorphy domain.
Uniform continuity makes $H_t$ nonzero on $\Gamma$ for sufficiently
small complex $t$. The argument principle keeps exactly $m$ zeros
inside, counted with their actual multiplicities. Their complete
test trace is
\[
 Z_\phi(t)=\frac1{2\pi i}\int_\Gamma
       \phi(x)\frac{\partial_xH_t(x_0+x)}{H_t(x_0+x)}\,dx
       =\sum_{\text{cluster}}\phi(x_i(t)).
 \tag{LT2}
\]
The contour integrand is holomorphic in $t$ on a neighbourhood
of the compact contour, so $Z_\phi$ is holomorphic. No derivatives
of individual colliding roots are used.

On an annulus about $\Gamma$, $H_t/H_0$ is uniformly close to $1$.
Its logarithm is the convergent branch equal to zero at $t=0$.
For $n\ge1$, differentiation of this branch and integration by
parts on $\Gamma$ give
\[
 \boxed{[t^n]Z_\phi(t)
   =-\operatorname{Res}_{x=0}
       \phi'(x)[t^n]\log(H_t/H_0)\,dx.}
 \tag{LT3}
\]
Here $[t^n]\log(H_t/H_0)$ is meromorphic inside $\Gamma$, with
its only possible pole at $0$, as the next finite formula proves.
The integration by parts follows by integrating the derivative
of the single-valued product $\phi[t^n]\log(H_t/H_0)$.
At order zero, $Z_\phi(0)=m\phi(0)$.

\paragraph{LT2. A closed finite formula.}
Let $E=x\,d/dx$, and set $(a)_{\underline r}
=a(a-1)\cdots(a-r+1)$, with empty product $1$.
For $a\ge1$ put
\[
 C_a[u](x)=\frac{(-1)^a}{a!}
             (E+m)_{\underline{2a}}u(x).
 \tag{LT4}
\]
The elementary identity
$\partial_x^{2a}(x^mu)=x^{m-2a}
(E+m)_{\underline{2a}}u$ follows by applying both sides to each
convergent monomial and differentiating the convergent series.
The PDE therefore determines its actual Taylor coefficient:
\[
 [t^a]H_t(x_0+x)=\frac{(-1)^a}{a!}
                       \partial_x^{2a}(x^mu),\qquad
 \frac{[t^a]H_t}{H_0}=x^{-2a}\frac{C_a[u]}u.
 \tag{LT5}
\]
Only finitely many time derivatives are used at each order.
In particular no convergence assertion about an unrestricted
formal heat operator is needed.

Define the following finite expression, holomorphic at $x=0$:
\[
 Q_{m,n}[u](x)=
 \sum_{\substack{k_1,\ldots,k_n\ge0\\
                  \sum_{a=1}^n a k_a=n}}
 \frac{(-1)^{K+1}(K-1)!}{\prod_{a=1}^n k_a!}
       \prod_{a=1}^n\left(\frac{C_a[u](x)}{u(x)}\right)^{k_a},
 \qquad K=\sum_{a=1}^n k_a.
 \tag{LT6}
\]
Expand the logarithm through time degree $n$, using
\[
 \log(1+Y)=\sum_{K=1}^n\frac{(-1)^{K+1}}K Y^K+O(Y^{n+1}),
 \qquad
 Y=\sum_{a=1}^n t^a x^{-2a}\frac{C_a[u]}u+O(t^{n+1}).
\]
The multinomial coefficient is $K!/\prod k_a!$; every selected
term has pole factor $x^{-2\sum a k_a}=x^{-2n}$.
This proves exactly
\[
 [t^n]\log(H_t/H_0)=x^{-2n}Q_{m,n}[u](x).
 \tag{LT7}
\]
In particular the pole order is at most $2n$.
Writing $\phi(x)=\sum_{k\ge0}\phi_kx^k$ in LT3 gives the
requested closed receiving formula
\[
 \boxed{[t^n]Z_\phi
   =-\sum_{k=1}^{2n}k\phi_k
                 [x^{2n-k}]Q_{m,n}[u].}
 \tag{LT8}
\]
Every sum and every coefficient extraction in LT4--8 is finite
at fixed $n$. All constants, signs and factorials refer to the
original $t$ and $x$.

\paragraph{LT3. The sufficient jet order, with universal boundaries.}
Introduce $\gamma_r=u_r/u_0$ solely as recorded coefficient ratios;
the function $H_t$ and its complete unit $u$ have not been rescaled.
The coefficient of $x^d$ in $C_a[u]/u$ is a polynomial in
$\gamma_1,\ldots,\gamma_d$, because the reciprocal series is
obtained recursively from $u\cdot u^{-1}=1$. It is homogeneous
of weight $d$ when $\gamma_r$ is given weight $r$:
the coefficient of $x^d$ in a product is a sum of products whose
indices add to $d$. LT6 preserves that property. Thus
$[x^d]Q_{m,n}[u]$ is such a polynomial of weight $d$.

For a general fixed test, LT8 depends only on $\phi_1,\ldots,\phi_{2n}$
and $u_0,\ldots,u_{2n-1}$, equivalently $u\bmod x^{2n}$.
The constant $\phi_0$ contributes only at time zero.
If $\phi$ vanishes to order at least $\ell\ge1$, then:
\[
 \begin{array}{c|c}
  \ell>2n &[t^n]Z_\phi=0\\
  \ell=2n &[t^n]Z_\phi=p_{2n}\phi_{2n}\\
  1\le\ell<2n &u_0,\ldots,u_D\text{ suffice},\quad D=2n-\ell.
 \end{array}
 \tag{LT9}
\]
Here $p_{2n}$ is the $2n$th root moment of the original-scale
heat polynomial
\[
 P_m(y)=\sum_{a=0}^{\lfloor m/2\rfloor}
       \frac{(-1)^am!}{a!(m-2a)!}y^{m-2a}.
 \tag{LT10}
\]
More generally $p_a$ denotes the sum of the $a$th powers of
all $m$ roots, counted with multiplicity, and $p_0=m$.
To prove the middle entry, LT8 uses only $Q_{m,n}[u](0)$,
which depends solely on $m$. At the auxiliary unit $u=1$,
$H_t=e^{-t\partial_x^2}x^m=t^{m/2}P_m(x/\sqrt t)$;
its roots are $\sqrt t$ times the roots of $P_m$.
Their $2n$th power sum is exactly $p_{2n}t^n$, proving that
constant. This is a calculation of a universal coefficient of the
retained family, not its replacement by a model at other orders.
The first entry follows because the sum in LT8 is empty.
If $\phi-\phi(0)$ has a higher order than $\phi$ itself, that
higher order may be used for every positive time coefficient.

\paragraph{LT4. Exact supply by the marked residue tower.}
For every $r\ge1$ retain
\[
 A_r=\mathbb C[x]/(x^r),\qquad
 \lambda_r(f)=\operatorname{Res}_0\frac{f(x)}{x^ru(x)}\,dx,
 \qquad B_r(f,g)=\lambda_r(fg).
 \tag{LT11}
\]
If $u^{-1}=\sum_{a\ge0}q_ax^a$, then
$\lambda_r(x^j)=q_{r-1-j}$ for $0\le j<r$.
The pairing matrix has entries $q_{r-1-i-j}$ when $i+j<r$,
and zero otherwise. Reversing columns gives a triangular matrix
with diagonal $q_0=u_0^{-1}$, so
\[
 \det B_r=(-1)^{r(r-1)/2}u_0^{-r}\ne0.
 \tag{LT12}
\]
Its entire marked pairing gives exactly $q_0,\ldots,q_{r-1}$,
since $\lambda_r(f)=B_r(1,f)$. The inverse recursion
\[
 u_0=q_0^{-1},\qquad
 u_a=-q_0^{-1}\sum_{b=1}^a q_bu_{a-b}\quad(1\le a<r)
 \tag{LT13}
\]
recovers exactly $u\bmod x^r$.
The original coordinate class $x$ and the given functional are
part of the data.

For $R\ge r$ the quotient algebra map $\pi_{R,r}:A_R\to A_r$
has kernel $(x^r)/(x^R)$. Its transpose under these perfect
pairings is the injective module map
\[
 j_{r,R}(f)=x^{R-r}\widetilde f,\qquad
 B_R(a,j_{r,R}f)=B_r(\pi_{R,r}a,f),\qquad
 \lambda_R(j_{r,R}f)=\lambda_r(f).
 \tag{LT14}
\]
A change of lift adds a multiple of $x^R$, proving
well-definedness. The coefficients of $x^{R-r},\ldots,x^{R-1}$
prove injectivity; cancelling $x^{R-r}$ in the residue proves
both identities. It is an $A_R$-module map when $A_r$ has its
action through $\pi$, not an algebra inclusion. Its image is
$(x^{R-r})=(\ker\pi)^{\perp_{B_R}}$, since the inclusion follows
from LT14 and both sides have dimension $r$. All compositions
are the literal quotient compositions and multiplication by the
corresponding powers of $x$.

Therefore $A_{2n}$ supplies the general LT8 coefficient, and
$A_{D+1}$ supplies the third case of LT9. When $m\ge D+1$,
the original multiplicity-$m$ receiver $A_m$ already supplies
these data by LT14. When $m<D+1$, the displayed larger object
is a specified analytic thickening. It does not add zero
multiplicity. The full inverse system determines all Taylor
coefficients of the original analytic $u$ by LT13, and hence
determines that germ by the uniqueness of a convergent power
series. A single finite receiver supplies only its stated jet.

On the full supported module
$\{(0,\lambda):\lambda\in L\}\cup\{(a,1_L):a\in A_r\}$,
each map retains the label:
$(a,1_L)\mapsto(\pi a,1_L)$ or $(ja,1_L)$ and
$(0,\lambda)\mapsto(0,\lambda)$. Addition and the prescribed
scalar action commute with these rules directly. The quotient
is multiplicative; $j$ has the module property just proved.
An amplitude in $\ker\pi$ becomes the top-supported zero, not
external absence. The trace receiver LT8 is a polynomial
function of the recorded unit data, not an asserted algebra map;
its output likewise retains its specified scalar support.

\paragraph{LT5. The highest-jet coefficient is exactly computable.}
Suppose $1\le\ell<2n$, $\phi_\ell\ne0$, and put $D=2n-\ell$.
Weight homogeneity implies that the dependence on $\gamma_D$
in LT8 is affine, and its coefficient is independent of every
lower unit coefficient. Only the summand $k=\ell$ can contain it.
The exact coefficient is
\[
 \boxed{\frac{\partial}{\partial\gamma_D}[t^n]Z_\phi
       =\phi_\ell\,\sigma_{m,n,\ell},\qquad
 \sigma_{m,n,\ell}
       =-\ell[y^{-\ell}]\frac{P_{m+D}(y)}{P_m(y)}.}
 \tag{LT15}
\]
The Laurent coefficient is at infinity. To prove it, it is
enough by the preceding weight argument to evaluate at
$u(x)=u_0(1+\eta x^D)$. The actual finite polynomial heat
solution with this initial value is
\[
 H^{(\eta)}(t,x)=u_0\left(
 e^{-t\partial_x^2}x^m+\eta e^{-t\partial_x^2}x^{m+D}\right).
 \tag{LT16}
\]
The exponentials here terminate, and direct differentiation proves
the heat equation and exact initial multiplicity $m$.
Differentiating the logarithm at $\eta=0$ gives
$H_{m+D}(t,x)/H_m(t,x)-x^D$; its second term has zero
coefficient at positive time degree. With
$H_k(t,x)=t^{k/2}P_k(x/\sqrt t)$, the required term is
$x^{-\ell}[y^{-\ell}]P_{m+D}(y)/P_m(y)$.
Substitution in LT3 proves LT15, including its minus sign.
The finite heat examples in LT16 test information content; no
new multiple zero of the arithmetic family is asserted.

There is an integer recurrence evaluating every multiplier without
approximating any root. Put
\[
 a_{k,j}=\frac{(-1)^j(k)_{\underline{2j}}}{j!},\qquad
 b_0=1,\quad
 b_j=a_{m+D,j}-\sum_{h=1}^j a_{m,h}b_{j-h}
       \quad(1\le j\le n).
 \tag{LT17}
\]
Terms with $2j>k$ are zero for the nonnegative integer $k$.
Dividing
$\sum_j a_{m+D,j}z^j$ by $\sum_j a_{m,j}z^j$ coefficientwise
proves $\sigma_{m,n,\ell}=-\ell b_n$.
Thus a zero multiplier is identified by exact finite arithmetic.
When it is zero, the highest coefficient is absent for every
unit. When it is nonzero, changing $\gamma_D$ by a nonzero
$\eta$ changes the trace coefficient by precisely
$\eta\phi_\ell\sigma_{m,n,\ell}$, while leaving the complete
$A_D$ receiver unchanged. This explicitly proves the necessity
of $A_{D+1}$ in that case.

\paragraph{LT6. A finite exact test for every remaining jet.}
The zero cases of LT15 require calculation, rather than a sharpness
claim based solely on the upper bound. The following formula computes
every lower sensitivity and thereby determines the exact cutoff.
In formal variables $z,x$, define
\[
 \begin{aligned}
 w(x)&=1+\sum_{r\ge1}\gamma_rx^r,\qquad\gamma_0=1,\\
 {\cal F}_m(z,x)&=
 \sum_{a=0}^{n}\sum_{r\ge0}\gamma_r a_{m+r,a}z^ax^r,
 \\
 V_{m,n,r}(x)&=[z^n]\frac{\sum_{a=0}^{n}a_{m+r,a}z^a}
                                      {{\cal F}_m(z,x)}.
 \end{aligned}
 \tag{LT18}
\]
Only coefficients of $x$ through $2n-\ell-r$ need be formed
below. All divisions are finite formal divisions with nonzero
constant term; $w$ is a record of ratios, not a change of the
retained function.
For every $1\le r\le D$,
\[
 \boxed{\frac{\partial}{\partial\gamma_r}[t^n]Z_\phi
 =-\sum_{k=\ell}^{2n-r} k\phi_k
              [x^{2n-r-k}]V_{m,n,r}(x).}
 \tag{LT19}
\]
Indeed differentiation of $H_t$ with respect to $\gamma_r$
is $u_0H_{m+r}$, by the linearity of LT5. The positive-time
part of the logarithmic derivative is therefore
$[t^n]u_0H_{m+r}/H_t$. Factoring $x^m$ and putting
$z=t/x^2$ gives $x^{r-2n}V_{m,n,r}(x)$. The derivative
of $\log H_0$ is time-independent, and LT3 proves LT19.

For each fixed $m,n$ and fixed test jet, the right sides of LT19
are explicitly determined finite polynomials in
$\gamma_1,\ldots,\gamma_{D-r}$. This list may include the
differentiated variable $\gamma_r$ when $2r\le D$.
Let $r_*$ be the greatest $r$ for which this polynomial
has a nonzero coefficient, with $r_*=0$ if none do. This is a
finite coefficient test, not an unresolved analytic condition.
LT8 is a polynomial, so it is independent of all
$\gamma_r$ with $r>r_*$, and $A_{r_*+1}$ supplies it.
If $r_*>0$, the nonzero polynomial in LT19 cannot vanish at
every tuple of real rational values of its unit variables: induction on
the number of variables reduces that assertion to the fact that
a nonzero one-variable polynomial has finitely many roots.
Choose such a rational basepoint in all variables occurring
in the derivative, then vary only $\gamma_{r_*}$ near that
basepoint. The derivative there is nonzero, so a sufficiently
small nonzero rational change gives a different value. The resulting
finite polynomial heat families have identical receiver $A_{r_*}$
and different trace coefficients. Thus the stated finite test
determines the exact least receiver order for the fixed test over
the class of genuine analytic heat families. If $r_*=0$, the
coefficient is universal and needs no unit data.

\paragraph{LT7. Explicit zero multipliers and their actual receivers.}
LT17 gives the following complete lists at the first three
positive time orders:
\[
\begin{array}{c|c|l}
n&\ell&\sigma_{m,n,\ell}\\ \hline
1&1&2m\\
2&1&-12m\\
2&2&8m(m-1)\\
2&3&12m(m-1)\\
3&1&-40m(m-4)\\
3&2&-128m(m-1)\\
3&3&24m(m-1)(m-5)\\
3&4&32m(m-1)(2m-3)\\
3&5&40m(m-1)(2m-3).
\end{array}
\tag{LT20}
\]
For example the first exceptional case $m=4,n=3,\phi=x$ has
the exact coefficient
\[
 [t^3]Z_x=-32\gamma_1^5+160\gamma_1^3\gamma_2
       -96\gamma_1^2\gamma_3-192\gamma_1\gamma_2^2
       +192\gamma_2\gamma_3.
 \tag{LT21}
\]
Both $\gamma_5$ and $\gamma_4$ are absent, and the actual least
receiver is $A_4$, since the $\gamma_3$ derivative
$192\gamma_2-96\gamma_1^2$ is a nonzero polynomial.
For a fixed test of exact order $1$ at the same $(m,n)$,
\[
 \partial_{\gamma_4}[t^3]Z_\phi=-1536\phi_2,\qquad
 \left.\partial_{\gamma_3}[t^3]Z_\phi\right|_{\phi_2=0}
       =\phi_1(192\gamma_2-96\gamma_1^2)-288\phi_3.
 \tag{LT22}
\]
Thus the least receiver is $A_5$ when $\phi_2\ne0$ and $A_4$
when $\phi_2=0$. In the latter case $\phi_1\ne0$ ensures that
the second polynomial is nonzero.

For $m=5,n=3,\phi=x^3$ the result is
\[
 [t^3]Z_{x^3}=-440\gamma_1^3+960\gamma_1\gamma_2.
 \tag{LT23}
\]
For a fixed test of exact order $3$, its highest remaining
derivative is $960\phi_3\gamma_1+4480\phi_4$, so the least
receiver is $A_3$. The term $\gamma_3$ is absent by LT20.

Zero multipliers occur beyond these first orders. An exact instance is
\[
 \sigma_{m,5,2}=-1024m(m-1)(m-4)(m-21).
 \tag{LT24}
\]
At $m=4$ and $m=21$, for any fixed test of exact order $2$,
the degree-$8$ unit coefficient is absent, while the next
derivatives are respectively
\[
 \begin{aligned}
 \partial_{\gamma_7}[t^5]Z_\phi
       &=32256\phi_2\gamma_1+532224\phi_3 &&(m=4),\\
 \partial_{\gamma_7}[t^5]Z_\phi
       &=-1270080\phi_2\gamma_1-182891520\phi_3 &&(m=21).
 \end{aligned}
 \tag{LT25}
\]
Both polynomials are nonzero since $\phi_2\ne0$, so $A_8$ is
the exact least receiver in these cases.
The finite derivation of LT20--25 is substitution of LT4 into
LT6, with the coefficient condition $\sum a k_a=n$ retained;
LT19 gives the displayed sensitivities directly. The included
exact checker evaluates both this partition calculation and an
independent local factor/Newton calculation. LT17 and LT19
retain all other exceptional cases, without a claim that this
displayed list exhausts larger time orders.

\paragraph{LT8. The complete simple-zero refinement.}
At $m=1$, the implicit function theorem gives a single
holomorphic zero $x(t)=\sum_{a\ge1}\alpha_at^a$.
Applied to the test $x$, LT9 shows that $\alpha_a$ uses
unit coefficients only through degree $2a-1$.
The coefficient of $\gamma_{2a-1}$ is, by LT15,
\[
 -[y^{-1}]\frac{P_{2a}(y)}{P_1(y)}
     =(-1)^{a+1}\frac{(2a)!}{a!},
 \qquad \alpha_1=2\gamma_1.
 \tag{LT26}
\]
Here $P_1=y$, and the constant term of $P_{2a}$ is
$(-1)^a(2a)!/a!$, proving the formula.
For a fixed test of exact order $\ell\ge1$, the coefficient
$[t^n]\phi(x(t))$ is zero if $n<\ell$. If $n\ge\ell$, a product
of $k$ root coefficients contributing to time degree $n$ has
each index at most $n-k+1\le n-\ell+1$. It therefore uses
unit coefficients only through
\[
 R=2n-2\ell+1.
 \tag{LT27}
\]
This is a stronger bound than LT9 for every $\ell>1$.
When $n>\ell$, the unique way to include
$\alpha_{n-\ell+1}$ at time degree $n$ is to take it once
and take $\alpha_1$ in all other $\ell-1$ factors.
Consequently the highest-unit derivative is
\[
 \partial_{\gamma_R}[t^n]Z_\phi
   =\ell\phi_\ell(2\gamma_1)^{\ell-1}
       (-1)^{n-\ell+2}
          \frac{(2n-2\ell+2)!}{(n-\ell+1)!}.
 \tag{LT28}
\]
For $n=\ell$, the coefficient is
$\phi_\ell(2\gamma_1)^\ell$, whose derivative is the same
formula with $R=1$. The polynomial in LT28 is nonzero, so
$A_{2n-2\ell+2}$ is the exact least receiver over analytic
simple-zero families. In particular LT15's zero multipliers
for $m=1,\ell>1$ are actual losses of highest-jet information,
not exceptions to be suppressed.

\paragraph{LT9. Return to every original fixed-test flag.}
For fixed arithmetic tests $F,G$, put
$A(s)=F^\#(s)G(s)$, where
$F^\#(s)=\overline{F(1-\bar s)}$, and
$\phi(x)=A(\rho+ix/2)$. Its coefficients are exactly
\[
 \phi_k=(i/2)^k\frac{A^{(k)}(\rho)}{k!}
       =\sum_{a+b=k}\overline{(i/2)^af_a}\,(i/2)^bg_b,
 \quad f_a=F^{(a)}(\rho)/a!,\quad g_b=G^{(b)}(\rho)/b!.
 \tag{LT29}
\]
Reflection fixes the centre and conjugates the coefficients in
the spatial coordinate, proving the second identity.
If both tests vanish to order at least $j$, then $\phi$ vanishes
to order at least $2j$. Thus at $n=j$ the coefficient is the
universal $p_{2j}\bar f_jg_j/4^j$. At $n=j+1$, $j\ge1$,
LT9 supplies the full coefficient from $A_3$; all actual test
jets through total degree $2j+2$ remain in LT8.
Writing $c=\gamma_1$, $d=\gamma_2$, and $b=2d-c^2$ gives
\[
\begin{split}
[t^{j+1}]Z_{F^\#G}=\frac1{4^j}\bigg[&
 \bigl(2jb\,p_{2j}+4j(2j-1)c^2p_{2j-2}\bigr)\bar f_jg_j\\
 &+i(2j+1)c\,p_{2j}
            (\bar f_jg_{j+1}-\bar f_{j+1}g_j)\\
 &+\frac{p_{2j+2}}4
       (\bar f_{j+1}g_{j+1}-\bar f_jg_{j+2}-\bar f_{j+2}g_j)
 \bigg].
\end{split}
\tag{LT30}
\]
To verify the reduction directly, the recurrences
$P_{m+1}=yP_m-2P_m'$ and
$P_{m+2}(r)=-2rP_m'(r)$ at each root of $P_m$ give
the first two perturbed root terms
$r\sqrt t+2ct+b r t^{3/2}$. Expanding powers gives
$[t^{j+1}]\sum x_i^{2j}
=2jb p_{2j}+4j(2j-1)c^2p_{2j-2}$,
$[t^{j+1}]\sum x_i^{2j+1}=2(2j+1)cp_{2j}$,
and $[t^{j+1}]\sum x_i^{2j+2}=p_{2j+2}$.
The Hermite roots are simple. Indeed expanding
$e^{yz-z^2}$ proves both
$e^{yz-z^2}=\sum_{m\ge0}P_m(y)z^m/m!$ and the exact identity
$P_m(y)e^{-y^2/4}=(-2)^m\partial_y^m e^{-y^2/4}$.
Integration by parts $m$ times, with Gaussian-decaying boundary
terms, makes $P_m$ orthogonal to every lower-degree polynomial
for the measure $e^{-y^2/4}dy$. If it had fewer than $m$ real
sign changes, the product of the factors at those sign changes
would have degree less than $m$ and its product with $P_m$
would have a fixed nonzero sign almost everywhere. Its integral
could not vanish, a contradiction. Thus all $m$ roots are real
and simple.
The implicit expansions therefore apply to each root, and their
symmetric sum is even in $\sqrt t$. Substitution of LT29 gives
LT30, including every affine factor. This is precisely HR4 and
RJ20, not a replacement of their full fixed-test jets.
At every later $n$, LT8 and LT19 provide the entire original
coefficient and its exact information cutoff. Neither the
multiplicity nor the heat time has changed.

\begin{figure}[p]
\centering
\begin{tikzpicture}[>=Stealth,box/.style={draw,rounded corners,
 align=center,text width=11cm,inner sep=8pt}]
\node[box] (res) at (0,0)
 {Marked residue data on $A_{D+1}$\\
 $\lambda_{D+1}(x^{D-r})=q_r$, $0\le r\le D$};
\node[box] (unit) at (0,-2.1)
 {Actual unit coefficients $u_0,\ldots,u_D$\\
 LT13; $D=2n-\ell$ for a test of order $\ell<2n$};
\node[box] (finite) at (0,-4.2)
 {Finite partition polynomial $Q_{m,n}[u]$\\
 $[t^n]Z_\phi=-\sum_{k=\ell}^{2n}k\phi_k[x^{2n-k}]Q_{m,n}[u]$};
\node[box] (sharp) at (0,-6.6)
 {Exact information test\\
 highest coefficient: $-\ell[y^{-\ell}]P_{m+D}/P_m$\\
 every lower coefficient: LT19; zero multipliers retained};
\draw[->] (res)--node[right]{inverse unit recursion}(unit);
\draw[->] (unit)--node[right]{LT4--8}(finite);
\draw[->] (finite)--node[right]{LT15--19}(sharp);
\end{tikzpicture}
\caption{The all-order receiving map in the original coordinate.
LT3 proves the contour identity; LT6--8 give the finite formula;
LT11--14 prove the marked residue maps; LT15--28 determine
sharpness and its exceptions. The actual multiplicity is still $m$.
The diagram's $A_{D+1}$ is an analytic data receiver, and the
original heat source and earlier fixed-test proofs are cited at
the start of this section.}
\end{figure}

\clearpage
\section{A certified arithmetic crossing with a surviving negative direction}
\label{sec:native-crossing}

The preceding WR1--14, reproduced in full in this edition, receive the
exact vector in the supplied \texttt{rational\_witness.json} into the
complete Rodgers--Tao theta trace. The companion manuscript
\emph{A non-Eulerian RH counterexample attempt}, equations (4.2)--(7.2),
supplied the vector and the signs at times $-2$ and $0$. Those signs
are credited to that source. The calculation here adds a uniform
curvature proof, an isolated simple crossing, and a second rational
test that is strictly negative at that crossing. It uses the original
heat time, coefficients, spatial coordinate and every root contribution.
The human analytic sources are Brad Rodgers and Terence Tao,
\href{https://arxiv.org/abs/1801.05914v5}{\emph{The de Bruijn--Newman
constant is non-negative}}, original author equations
\texttt{htdef}, \texttt{phidef}, \texttt{hoz}, \texttt{sas}, and
Jacques Hadamard's
\href{https://www.numdam.org/item/JMPA_1893_4_9__171_0/}{1893
entire-function factorization paper}; WR1--5 articulate their precise
use and the original-coordinate receiving maps.

\paragraph{NC1. The fixed original tests.}
Retain exactly
\[
 f_*(Z)=\sum_{i=0}^{18}c_i Z^{-2i-1},\qquad
 c_i=1000^i n_i/10^{120},\qquad g_0(Z)=Z^{-1},
 \tag{NC1}
\]
where the nineteen integers $n_i$ are the complete unrounded vector
in the retained JSON. Its SHA256 is
\path{85314e9fb3115647b7f03a5e00f1b5512c4b23c9280a753f3cf2cababdfed4ac}.
The functions computed below are
\[
 \begin{split}
 Q(t)&=Q_t(f_*,f_*)=\sum_{i,j=0}^{18}c_ic_j S_{i+j+1}(t),\\
 B(t)&=Q_t(g_0,f_*)=\sum_{i=0}^{18}c_i S_{i+1}(t),\\
 A(t)&=Q_t(g_0,g_0)=S_1(t).
 \end{split}\tag{NC2}
\]
All $c_i$ are real. WR1--4 prove the complete-root interpretation,
absolute convergence, and real analyticity on every compact real
time interval. In particular $H_t(0)>0$, so these tests have no pole
at a zero of $H_t$. From WR4 at $n=1$ one obtains the additional
strict identity $A(t)=M_1(t)/M_0(t)>0$.

\paragraph{NC2. Every required derivative from the original moments.}
Put $D_{r,n}(t)=\partial_t^r S_n(t)$ for $r\ge0$, $n\ge1$.
Differentiating WR5 $r$ times by the product rule gives
\[
 D_{r+1,n}=2n\left[
 \sum_{a=1}^n\sum_{j=0}^r\binom rj
 D_{j,a}D_{r-j,n+1-a}-(2n+1)D_{r,n+1}\right].
 \tag{NC3}
\]
This is an identity of analytic functions, so it remains valid
through collisions away from the fixed pole. Its proof uses no
division by root differences. The entire initial row $D_{0,n}$
is computed from WR4, including the positive denominator $M_0$.
Induction on $r$ shows that derivatives through order $r$ of $Q$
use precisely moments through $37+r$; the corresponding upper
bound for $B$ is $19+r$. Hence moments through $42$ supply every
entry used here, including fifth derivatives.

There is an independent formal verification of the differential
identity. In a variable $h$, set
\[
 a_n(h)=\frac{(-1)^n}{(2n)!}
       \sum_{r\ge0}\frac{M_{n+r}(t)}{r!}h^r.
 \tag{NC4}
\]
The identity $M_n^{(r)}=M_{n+r}$ follows by differentiating the
dominated integral in WR1. Divide the $Z$-series $\partial_ZH/H$
with coefficients $a_n(h)$. Equating the resulting coefficients
gives the same $D_{r,n}/r!$. The included exact checker compares
NC3 with this separate series-division construction; the analytic
proof remains WR4--5 and the product rule above.

\paragraph{NC3. Uniform interval calculation on the full time window.}
Set $h_0=10^{-20}$ and
\[
 I=[-2,-2+h_0].\tag{NC5}
\]
The script \texttt{native\_curvature.py} computes all moments
$M_0,\ldots,M_{42}$ both at $t=-2$ and uniformly for $t\in I$.
It uses the same $160$ bracketed Gauss--Legendre nodes per panel,
$24$ panels on $[0,3]$, first $16$ theta summands and directed
$360$-digit intervals as WR8--11. The node brackets and positive
weights are certified again by their exact recurrence signs;
decimal Newton iterates only propose those brackets.

Here the panel estimate holds uniformly, not merely at two times.
For real $t\in I$ and complex $z=x+iy$ on the ellipse,
\[
 |e^{tz^2}|=e^{t(x^2-y^2)}\le e^{2B^2},\qquad B=15/128.
 \tag{NC6}
\]
Thus the full WR9 panel error is bounded by
$(4\ell/3)P_v e^{2B^2}R^{2k}4^{-319}$ for every $0\le k\le42$.
The omitted theta summands on $[0,3]$ still have error at most
$192\,17^4\,3^{2k}e^{-840}$, since $t\le0$.
The extended spatial-tail proof, for $u\ge3$, is
\[
 2k\log u+9u\le93u\le31u^2\le e^{4u}.
 \tag{NC7}
\]
The last inequality holds at $3$, and $e^{4u}/u^2$ is increasing
thereafter. The exact same bound
$64\int_3^\infty e^{-2e^{4u}}du<e^{-2000}$ therefore applies.
These estimates prove uniform containment of each actual moment.
Outward arithmetic in WR4 and NC3 then proves uniform containment
of every derivative in the retained receipt.

In particular the uniform fifth-derivative bound is
\[
 \sup_{t\in I}|Q^{(5)}(t)|<4.8\,10^{17}=:C_5.
 \tag{NC8}
\]
Direct interval evaluation of $Q$ on $I$ has a wide enclosure
because it repeats the same time variable in many cancelling terms.
NC8 retains a valid consequence of that enclosure. Taylor's theorem
uses it together with the accurately evaluated derivatives at the
original endpoint; no cancellation is assumed or discarded.
The endpoint values, with the complete binary enclosures in
\texttt{NATIVE\_CURVATURE\_INTERVALS.json}, begin
\[
\begin{array}{c|r}
r&Q^{(r)}(-2)\\\hline
0&-2.3840641788179150923540341723\ldots\,10^{-46}\\
1& 1.0766217789670091488291342795\ldots\,10^{-44}\\
2& 0.0000142941887834022346441600660202\ldots\\
3& 0.0000713941668222182582756860423090\ldots\\
4& 0.0002385266580194101433481063682461\ldots
\end{array}\tag{NC9}
\]
These displayed decimals identify the values; the calculation uses
their exact outward binary endpoints, not the truncated decimals.

\paragraph{NC4. Positive curvature and the unique simple crossing.}
For $0\le d\le h_0$ Taylor's theorem with integral remainder gives,
for $r=0,1,2$,
\[
 \left|Q^{(r)}(-2+d)-
 \sum_{j=r}^4\frac{Q^{(j)}(-2)}{(j-r)!}d^{j-r}\right|
 \le\frac{C_5d^{5-r}}{(5-r)!}.
 \tag{NC10}
\]
Indeed the integral remainder is
$\int_0^d Q^{(5)}(-2+y)(d-y)^{4-r}dy/(4-r)!$.
Applying the interval endpoint values in NC9 proves on all of $I$
\[
 0.00001429418878340223<Q''(t)
 <0.00001429418878340224.
 \tag{NC11}
\]
Since $Q'(-2)>0$, $Q$ is strictly increasing on $I$.
The same NC10 enclosure gives $Q(-2)<0$ and $Q(-2+h_0)>0$.
The intermediate value theorem and strict monotonicity prove exactly
one zero, which we denote $t_*=-2+d_*$. Its derivative is positive.

The script \texttt{certify\_native\_crossing.py} performs $120$
rational bisections, using NC10 at each proposed rational endpoint.
Every sign must be strict; an undecided interval stops the code.
The resulting certified bracket is
\[
 \frac{767703802914770771337427964528646841}
      {2^{120}10^{20}}
 <d_*<
 \frac{767703802914770771337427964528646842}
      {2^{120}10^{20}}.
 \tag{NC12}
\]
It implies the more readable bounds
\[
\begin{aligned}
d_*&>5.775561493959038849338863555346879146\,10^{-21},\\
d_*&<5.775561493959038849338863555346879155\,10^{-21}.
\end{aligned}\tag{NC13}
\]
The derivative obeys
\[
 \begin{aligned}
 Q'(t_*)&>8.25569663247991463159849730318\,10^{-26},\\
 Q'(t_*)&<8.25569663247991463159849730320\,10^{-26}.
 \end{aligned}
 \tag{NC14}
\]
In particular this is a transversal crossing of the scalar test.
The positive root $d_2$ of the endpoint quadratic is also evaluated
without changing the function:
\[
 d_2=\frac{-Q'(-2)+
 \sqrt{Q'(-2)^2-2Q''(-2)Q(-2)}}{Q''(-2)}.
 \tag{NC15}
\]
The same outward calculations give
\[
 -2.7767730799140320\,10^{-41}<d_*-d_2
 <-2.7767730799140312\,10^{-41}.
 \tag{NC16}
\]
This evaluates a nonlinear correction to the actual crossing;
the quadratic expression alone was not used to certify its value.

\paragraph{NC5. The complete two-test receiver at the crossing.}
Apply NC3 to $B$ in NC2. Its first derivative at $-2$ begins
$0.000128434311029871499097473571169\ldots$, while
$B(-2)$ lies near $1.81298975228\,10^{-123}$.
The full interval calculation bounds $|B^{(5)}|<0.063$ on $I$.
The analogous fourth-order Taylor enclosure at the entire interval
NC12 gives
\[
\begin{gathered}
\begin{aligned}
B(t_*)&>7.41780261287284496794922217392\,10^{-25},\\
B(t_*)&<7.41780261287284496794922217394\,10^{-25},
\end{aligned}\\
0.01109183050805782326708<A(t_*)<0.01109183050805782326710.
\end{gathered}\tag{NC17}
\]
The script retains the unrounded enclosures and checks $B(t_*)>0$.
On the original two-dimensional subspace with ordered basis
$(f_*,g_0)$, the actual trace form is exactly
\[
 \begin{pmatrix}0&B(t_*)\\B(t_*)&A(t_*)\end{pmatrix},
 \qquad\det=-B(t_*)^2<0.
 \tag{NC18}
\]
The two tests are linearly independent because $c_{18}=10^{54}$
and $g_0$ has no $Z^{-37}$ term. The real symmetric matrix therefore
has one positive and one negative eigenvalue. The vector $f_*$ is
isotropic but outside the radical: its pairing with $g_0$ is nonzero.
This statement is about the original full-root form, without deleting
the roots outside any selected cluster.

There is an explicit rational negative vector at the same time:
\[
 \widetilde f=f_*-10^{-22}g_0,
 \qquad
 Q_{t_*}(\widetilde f,\widetilde f)
 =-2\,10^{-22}B(t_*)+10^{-44}A(t_*).
 \tag{NC19}
\]
Substitution of the certified NC17 intervals proves
\[
 -3.74377471768786666882\,10^{-47}
 <Q_{t_*}(\widetilde f,\widetilde f)
 <-3.74377471768786666880\,10^{-47}.
 \tag{NC20}
\]
Thus the zero of this one fixed test is not a positivity boundary
for its two-test subspace. NC18--20 give the exact inclusion and
the surviving negative direction establishing that assertion.
Continuity also preserves that strict negative value on some open
interval about $t_*$; no size is claimed without a further estimate.

\paragraph{NC6. Supported zero, ordinary amplitude, and the retained label.}
Let $V=\operatorname{span}_{\mathbb C}\{f_*,g_0\}$ and retain the
original bounded distributive support lattice $L$. On the full
carrier $M_L(V)=(V\times\{1_L\})\cup(\{0\}\times L)$, use amplitude
addition with support join and scalar multiplication with support
meet, exactly as in BT22. Its zero is $\tau=(0,0_L)$; the top
supported zero is $e=(0,1_L)$. The complete lift is the explicit map
\[
 \widehat Q_t((v,\lambda),(w,\mu))
       =(Q_t(v,w),\lambda\wedge\mu).
 \tag{NC21}
\]
If a support is below $1_L$, the corresponding amplitude is zero,
so the right side belongs to $M_L(\mathbb C)$. Sesquilinearity
of WR3 verifies the amplitudes; distributivity of $L$ verifies
the additive support identities, and complex conjugation retains
the support. Thus NC21 is a sesquilinear map of the specified
supported modules. The amplitude projections commute with NC21
and $Q_t$ by their displayed formulas.

At $t_*$, NC21 sends the pair of supported copies of $f_*$ to $e$,
while its pairing with the supported $g_0$ has the positive amplitude
in NC17. Applying the amplitude projection to the retained
$2\times2$ array gives exactly NC18. The entire carrier still retains
all other support labels. The diagonal supported zero therefore
leaves $f_*$ present and outside the radical: NC17 is the explicit
nonzero cross-pairing detecting the same vector.

The physical time of this computed example is the negative number
in NC12. WR7 separately gives the positive value of $Q_*(0)$.
The exact map transporting its derivatives between those times is
the original recurrence NC3, not a relabelling of $t_*$. The
calculation supplies an arithmetic test-space crossing and its
support-sensitive receiver. It does not assert a counterexample
at the arithmetic endpoint $t=0$ or identify $t_*$ as a zero of
$H_t$ itself.

\paragraph{NC7. Verification scope.}
The numerical proof uses outward \texttt{mpmath.iv} operations,
explicit analytic quadrature and tail errors, positive-denominator
checks, and exact rational bisection. It is computer-assisted;
formal verification of the interval library is not claimed.
\texttt{NATIVE\_CURVATURE\_INTERVALS.json} retains exact binary
endpoints for the point and uniform derivatives.
\texttt{CERTIFIED\_NATIVE\_CROSSING.json} retains the bracket,
endpoint signs, curvature, derivative, cross-pairing and negative
companion test. The symbolic checker verifies an independent
finite consequence of NC3--4; it does not replace the analytic
proof or the interval calculation.

\begin{figure}[p]
\centering
\begin{tikzpicture}[>=Stealth,x=0.72cm,y=0.62cm]
\draw[->] (0,-2.8)--(0,5.8) node[above] {$10^{46}Q(-2+d)$};
\draw[->] (0,0)--(10.7,0) node[right] {$d/10^{-21}$};
\foreach \x in {2,4,6,8,10}{\draw(\x,0.08)--(\x,-0.08) node[below] {\small\x};}
\foreach \y in {-2,2,4}{\draw(0.08,\y)--(-0.08,\y) node[left] {\small\y};}
\draw[thick,blue,domain=0:10,samples=100]
 plot(\x,{-2.384064178817915+0.07147094391701117*\x*\x});
\draw[dashed] (5.775561494,-2.8)--(5.775561494,3.2);
\fill (5.775561494,0) circle(2pt);
\node[anchor=west,align=left] at(6.0,2.1)
 {$d_*/10^{-21}=5.77556149\ldots$\\unique simple crossing};
\node[anchor=west,align=left,text width=10cm] at(0,-4.5)
 {Curve: endpoint quadratic with the axis units shown.\\
 The quadratic expression's total error is below $10^{-14}$ in these units.};
\node[draw,rounded corners,align=center,text width=11.5cm,inner sep=8pt]
 at(5,-8.5)
 {At the actual crossing, in the original tests $(f_*,Z^{-1})$:\\[3pt]
 $\begin{pmatrix}0&B_*\\B_*&A_*\end{pmatrix}$,
 $B_*>0$, $A_*>0$, $\det=-B_*^2<0$\\[3pt]
 $f_*-10^{-22}Z^{-1}$ has value in
 $(-3.744,-3.743)\,10^{-47}$.};
\end{tikzpicture}
\caption{The actual fixed-witness crossing and its surviving negative
direction. The axes record explicit display units only; time remains
$t=-2+d$ and the original test coefficients are unchanged. The curve
uses the endpoint quadratic (NC9--10); the displayed error bound also
includes its omitted linear term and coefficient rounding. NC12--14
give the rigorous crossing, and NC17--20 give the complete two-test
matrix and negative companion. The full theta source is Rodgers--Tao
(cited above); the vector is the supplied rational witness.}
\end{figure}

\clearpage
\endgroup
